\documentclass[10pt,a4paper]{amsart}
\usepackage{amsmath,amssymb,amsthm,graphicx}
\usepackage{verbatim}
\newenvironment{reference}{\comment}{\endcomment}
%\newenvironment{reference}{}{}
\newenvironment{slogan}{\comment}{\endcomment}
\newenvironment{history}{\comment}{\endcomment}

% For commutative diagrams we use Xy-pic
\usepackage[all]{xy}

% We use 2cell for 2-commutative diagrams.
\xyoption{2cell}
% Theorem environments.
%
\theoremstyle{plain}
\newtheorem{theorem}[subsection]{Theorem}
\newtheorem{proposition}[subsection]{Proposition}
\newtheorem{lemma}[subsection]{Lemma}

\theoremstyle{definition}
\newtheorem{definition}[subsection]{Definition}
\newtheorem{example}[subsection]{Example}
\newtheorem{exercise}[subsection]{Exercise}
\newtheorem{situation}[subsection]{Situation}

\theoremstyle{remark}
\newtheorem{remark}[subsection]{Remark}
\newtheorem{remarks}[subsection]{Remarks}

\numberwithin{equation}{subsection}

% Macros
%
\def\lim{\mathop{\mathrm{lim}}\nolimits}
\def\colim{\mathop{\mathrm{colim}}\nolimits}
\def\Spec{\mathop{\mathrm{Spec}}}
\def\Hom{\mathop{\mathrm{Hom}}\nolimits}
\def\Ext{\mathop{\mathrm{Ext}}\nolimits}
\def\SheafHom{\mathop{\mathcal{H}\!\mathit{om}}\nolimits}
\def\SheafExt{\mathop{\mathcal{E}\!\mathit{xt}}\nolimits}
\def\Sch{\mathit{Sch}}
\def\Mor{\mathop{\mathrm{Mor}}\nolimits}
\def\Ob{\mathop{\mathrm{Ob}}\nolimits}
\def\Sh{\mathop{\mathit{Sh}}\nolimits}
\def\NL{\mathop{N\!L}\nolimits}
\def\CH{\mathop{\mathrm{CH}}\nolimits}
\def\proetale{{pro\text{-}\acute{e}tale}}
\def\etale{{\acute{e}tale}}
\def\QCoh{\mathit{QCoh}}
\def\Ker{\mathop{\mathrm{Ker}}}
\def\Im{\mathop{\mathrm{Im}}}
\def\Coker{\mathop{\mathrm{Coker}}}
\def\Coim{\mathop{\mathrm{Coim}}}

% Boxtimes
%
\DeclareMathSymbol{\boxtimes}{\mathbin}{AMSa}{"02}

%
% Macros for moduli stacks/spaces
%
\def\QCohstack{\mathcal{QC}\!\mathit{oh}}
\def\Cohstack{\mathcal{C}\!\mathit{oh}}
\def\Spacesstack{\mathcal{S}\!\mathit{paces}}
\def\Quotfunctor{\mathrm{Quot}}
\def\Hilbfunctor{\mathrm{Hilb}}
\def\Curvesstack{\mathcal{C}\!\mathit{urves}}
\def\Polarizedstack{\mathcal{P}\!\mathit{olarized}}
\def\Complexesstack{\mathcal{C}\!\mathit{omplexes}}
% \Pic is the operator that assigns to X its picard group, usage \Pic(X)
% \Picardstack_{X/B} denotes the Picard stack of X over B
% \Picardfunctor_{X/B} denotes the Picard functor of X over B
\def\Pic{\mathop{\mathrm{Pic}}\nolimits}
\def\Picardstack{\mathcal{P}\!\mathit{ic}}
\def\Picardfunctor{\mathrm{Pic}}
\def\Deformationcategory{\mathcal{D}\!\mathit{ef}}
\usepackage{cleveref}
\crefname{originallemma}{original lemma}{original lemmas}
\Crefname{originallemma}{Original lemma}{Original lemmas}
\makeatletter
\newlabel{lemma-Mittag-Leffler}{{02N1}{source}{}{sourcepointer-1}{}}
\newlabel{lemma-Mittag-Leffler@cref}{{[originallemma][0][]02N1}{[1][0][]source}{}{}{}}
\makeatother
\crefname{originallemma}{original lemma}{original lemmas}
\Crefname{originallemma}{Original lemma}{Original lemmas}
\makeatletter
\newlabel{lemma-exact-sequence-ML}{{070D}{source}{}{sourcepointer-2}{}}
\newlabel{lemma-exact-sequence-ML@cref}{{[originallemma][0][]070D}{[1][0][]source}{}{}{}}
\makeatother
\title{A limit of cohomology of an inverse system of complexes}
\author{The Stacks Project Authors}
\date{stacks-project:homology.tex:local-snapshot-20261012; GNU FDL 1.2 or later}
\begin{document}
\maketitle
\noindent\textit{Editorial note.} All delivered mathematical text is the author's own, copied from the pinned Stacks Project chapter file homology.tex; no statement or proof step was rewritten, reordered or removed.\par\medskip
\noindent\textit{Editorial note.} The proof is the author's complete argument; the closing sentence 'Putting everything together we win.' follows the three displayed exact sequences whose conclusions were derived just before it, so it closes the argument rather than deferring a step.\par\medskip
\noindent\textit{Editorial note.} Every labelled reference occurring inside the counted spans is carried in assembly.external\_references as a non-load-bearing navigation pointer (same-chapter ones at their definition site, cross-chapter ones at the invocation inside the counted span), so no load-bearing mathematics is replaced by a stub. No bibliography citation occurs inside the counted spans.\par\medskip
\noindent\textit{Editorial note.} The same-name check for this card scanned every staged Stacks chapter of this repository -- all sources/stacks-*/pinned-original-source.tex files, the card's own chapter excluded -- and the labels found are recorded in dedup.same\_name\_cross\_chapter on the card.\par\medskip
\noindent\textit{Editorial note.} The preamble is the chapter's own native preamble excerpt already used for this source: the comment-type reference, slogan and history environments (native preamble lines 8-13), the Xy-pic setup (native preamble lines 14-18) and the native theorem declarations with the mathematical macros (native preamble lines 154-217, the actual last line of the file).\par\medskip
\noindent\textit{Allowed standard prerequisites (editorial).}\begin{itemize}
\item The Mittag-Leffler condition and its behaviour under quotients. Same chapter, Lemma lemma-Mittag-Leffler. Supplies that a quotient system of a Mittag-Leffler system is Mittag-Leffler and the exactness statement for limits, applied three times in the course of the proof.
\item Mittag-Leffler along a short exact sequence with essentially constant quotient. Same chapter, Lemma lemma-exact-sequence-ML. Supplies the Mittag-Leffler property for the system of cocycles, deduced from the Mittag-Leffler sub-system and the essentially constant quotient system that the hypothesis provides.
\item Limits commute with kernels. Same chapter, the standard identification of the limit of the kernels with the kernel of the limit. Supplies the description of the limit of the cocycle systems as the cocycles of the limit complex.
\end{itemize}\medskip
\ifdefined\hypertarget\hypertarget{sourcepointer-1}{}\fi
\noindent\textit{Source reference.} Original lemma 02N1: sources/stacks-homology/pinned-original-source.tex, source lines 7890--7890. Navigation pointer invoked at proof line 8048; the counted span contains the pointer only and the cited result is not reproduced inside it, so the pointer is registered as a non-load-bearing navigation stub in the established form for Stacks material. No bibliography citation is involved.\par\medskip
\ifdefined\hypertarget\hypertarget{sourcepointer-2}{}\fi
\noindent\textit{Source reference.} Original lemma 070D: sources/stacks-homology/pinned-original-source.tex, source lines 7978--7978. Navigation pointer invoked at proof line 8056; the counted span contains the pointer only and the cited result is not reproduced inside it, so the pointer is registered as a non-load-bearing navigation stub in the established form for Stacks material. No bibliography citation is involved.\par\medskip
\begin{lemma}
\label{lemma-apply-Mittag-Leffler-again}
Let
$$
(A^{-2}_i \to A^{-1}_i \to A^0_i \to A^1_i)
$$
be an inverse system of complexes of abelian groups and denote
$A^{-2} \to A^{-1} \to A^0 \to A^1$ its limit. Denote
$(H_i^{-1})$, $(H_i^0)$ the inverse systems of cohomologies, and
denote $H^{-1}$, $H^0$ the cohomologies of $A^{-2} \to A^{-1} \to A^0 \to A^1$.
If $(A^{-2}_i)$ and $(A^{-1}_i)$ are ML and
$(H^{-1}_i)$ is essentially constant, then
$H^0 = \lim H_i^0$.
\end{lemma}

\begin{proof}
Let $Z^j_i = \Ker(A^j_i \to A^{j + 1}_i)$ and
$I^j_i = \Im(A^{j - 1}_i \to A^j_i)$.
Note that $\lim Z^0_i = \Ker(\lim A^0_i \to \lim A^1_i)$ as
taking kernels commutes with limits.
The systems $(I^{-1}_i)$ and $(I^0_i)$ have ML as quotients of
the systems $(A^{-2}_i)$ and $(A^{-1}_i)$, see
Lemma \ref{lemma-Mittag-Leffler}.
Thus an exact sequence
$$
0 \to (I^{-1}_i) \to (Z^{-1}_i) \to (H^{-1}_i) \to 0
$$
of inverse systems where $(I^{-1}_i)$ has ML
and where $(H^{-1}_i)$ is essentially constant by assumption.
Hence $(Z^{-1}_i)$ has ML by
Lemma \ref{lemma-exact-sequence-ML}.
The exact sequence
$$
0 \to (Z^{-1}_i) \to (A^{-1}_i) \to (I^0_i) \to 0
$$
and an application of
Lemma \ref{lemma-Mittag-Leffler}
shows that $\lim A^{-1}_i \to \lim I^0_i$ is surjective.
Finally, the exact sequence
$$
0 \to (I^0_i) \to (Z^0_i) \to (H^0_i) \to 0
$$
and
Lemma \ref{lemma-Mittag-Leffler}
show that $\lim I^0_i \to \lim Z^0_i \to \lim H^0_i \to 0$
is exact. Putting everything together we win.
\end{proof}

\end{document}
